\section{Motivation}
\label{sec:motivation}

\subsection{DSNs and Their Data Placement}

Decentralized storage networks (DSNs) provide blob storage service backed by decentralized storage resources collectively provided by permissionless nodes.
Some nodes are rational and aim to generate profit with their storage supply, while the other nodes are Byzantine, behaving arbitrarily and potentially maliciously.
DSNs can provide reliable storage service, ensuring the stored data can be successfully retrieved, given that there are at most a certain fraction (e.g. \nicefrac{1}{3}) of nodes are Byzantine.
They guarantee reliable storage by driving rational nodes to correctly store and serve the data with incentive and tolerating Byzantine nodes.

In this work, we focus on improving the data placement strategy in DSNs.
In storage systems, data placement means given a data block, how to determine which nodes should store redundancy for this data block and what kind of redundancy they store.
These storing nodes are called the placement group of the data block.
Storing data redundantly on multiple nodes is crucial for tolerating simultaneous failures across them.
For example, by storing $r$ copies of data on distinct nodes, the data will still be available even when $r-1$ of the nodes fail at the same time.
So, as long as the storage system can repair the redundancy before all $r$ nodes fail, the storage will be durable.

In practice, storing large number of data copies incur excessive storage overhead.
Storage systems commonly employ erasure coding to achieve better storage efficiency.
They encode a data block into $n$ fragments, any $k$ of which can recover the data block, and store them on $n$ nodes.
Then, the storage systems can tolerate simultaneous failure of any combination of $n-k$ nodes while incurring only $\nicefrac{n}{k}$ times of the data block storage on these nodes.
The rule of thumb is that with higher failure rate of the nodes (so they fail more frequently), the storage systems must be able to tolerate simultaneous failures of larger scale, and thus configure with larger placement groups.

Compared to traditional storage systems, one unique challenges of data placement in DSNs stem from the Byzantine nodes and limited governance infrastructure.
The DSNs who are configured with placement groups of $n$ nodes and can tolerate $k$ simultaneous failures of them must ensure that, no more than $f$ Byzantine nodes are selected into the $n$ nodes so that $f+f'>k$, where at most $f'$ benign simultaneous failures are expected from the $n-f$ rational nodes.
Otherwise, those Byzantine nodes may collude with each other and create simultaneous failures that lead to data loss.
The strict security model also forbid DSNs to leverage a trusted single party to manage data placement, but rely on blockchain to reach consensus on the governing routine to tolerate Byzantine nodes.
Blockchains are notorious for high cost of on-chain state and high transaction cost.
In DSNs, every blockchain node is effectively a full ledge storage controller and replicate maintenance effort of all data blocks.
Keeping low on-chain overhead is crucial for DSNs to be scalable and manage large volume of data.

In summary, we want to achieve the following goals of DSNs with their data placement strategies:
\begin{itemize}
    \item Tolerate up to a certain fraction of Byzantine nodes;
    \item Keep low and scalable storage overhead: the storage volume should be proportional to the total storage supply of the nodes;
    \item Minimize the on-chain state footprint.
\end{itemize}

\subsection{Previous DSNs and Their Limitations}
\label{sec:motivation:review}

\begin{table}[t]
\centering
% \renewcommand{\arraystretch}{1.5}
\begin{minipage}{\columnwidth}
\begin{tabularx}{\textwidth}{>{\raggedright\arraybackslash}X c c c}
\toprule
Protocol                       & Volume                             & \shortstack{On-chain\\Footprint} & \shortstack{Number of\\Faulty Nodes} \\
\midrule
Filecoin, Sia                  & \ProsCell $O(N\cdot S)$            & \ConsCell $O(D\cdot P)$   & \ConsCell $O(P)$          \\
Arweave, Permacoin, Retricoin  & \ConsCell $O(\frac{N\cdot S}{P})$  & \ProsCell $O(1)$          & \ConsCell $O(P)$          \\
Walrus                         & \ConsCell $O(S)$                   & \ProsCell $O(D)$          & \ConsCell $O(1)$          \\
Swarm                          & \ConsCell $O(\frac{N\cdot S}{P})$  & \ProsCell $O(1)$          & \ProsCell $O(N)$          \\
\sys                           & \ProsCell $O(N\cdot S)$            & \ProsCell $O(D)$          & \ProsCell $O(N)$          \\
\bottomrule
\end{tabularx}
\end{minipage}
\caption{
    Comparison of DSNs.
    $N$ indicates the total number of nodes, $S$ indicates the storage supply of each node (assuming uniform nodes), $D$ indicates the number of stored data, $P$ indicates placement group size.
}
\label{tab:dsn}
\end{table}

We summarize previous DSNs in \autoref{tab:dsn}, comparing them with three dimensions: storage volume, on-chain footprint and fault tolerance.
The storage volume and fault tolerance are the higher is better while the on-chain footprint is the lower is better.
Next, we will explain the data placement strategies employed by these DSNs and how they have caused various drawbacks in the DSNs.

\paragraph{Explicitly store placement on-chain.}

\paragraph{Implicit incentive-driven placement.}

\paragraph{Static placement.}

\paragraph{Randomized replication.}
